% Copia de 3A22/graficos/gfpu.tex (3.A.22, I.2); mismo código y cifras.
% Construcción de la gran frontera de posibilidades de utilidad (3.A.22, I.2).
% Izquierda: FPP X²/49 + Y²/36 = 1. E = (X*, Y*) = (5; 4,199); RMT_XY en E = 0,8748
%   (recta tangente: corta el eje X en 9,8). Caja del consumo con O_A = (0,0) y O_B = E,
%   U_A = X_A^0,4 Y_A^0,6 y U_B = X_B^0,6 Y_B^0,4 (las de la caja del consumo).
%   Curva de contrato: Y_A = 1,5 X_A Y* / ((2/3)(X* - X_A) + 1,5 X_A).
%   E* = (1,76; 2,3095): RMS^A = RMS^B = 0,8748 = RMT; U_A = 2,0717, U_B = 2,6114.
%   E_1 = (3; 5,421) y E_3 = (6,3; 2,615), otros 2 puntos de la FPP (solo marcados).
% Derecha (esquemático): ni las FPU ni la GFPU tienen por qué ser cóncavas (utilidad ordinal).
%   FPU_k: superelipses (U_A/a)^p + (U_B/b)^p = 1: FPU_1 (4; 7,5; 4,2), FPU_2 (6,2; 5,9; 4,2),
%   FPU_3 (7,6; 3,9; 3,5). GFPU (en polares): máximo suavizado de las 3 (suavizado cuadrático,
%   ε = 0,45) + 0,25·Π_k (1 - exp(-((θ - θ_k)/10°)²)), con θ_k = 69°, 45° y 19°: toca cada FPU
%   solo en E*_k (misma tangente) y queda por fuera en el resto (comprobado numéricamente).
%   Tramos cóncavos en torno a cada E*_k y 2 hondonadas convexas donde se cruzan las FPU
%   (cerca de (4,3; 6,4) y de (6,7; 3,4)); decreciente, de (0; 7,747) a (7,843; 0).
%   E*_1 = (2,729; 7,110), E*_2 = (5,121; 5,121), E*_3 = (7,135; 2,457).
% Autor: Víctor Gutiérrez Marcos
\documentclass[tikz,border=4pt]{standalone}
% ==========================================================================
% tikz-tcee.tex — Estilos comunes de los gráficos TikZ del temario
% Autor: Víctor Gutiérrez Marcos
%
% Cada gráfico es un fichero standalone en fuentes/<tema>/graficos/:
%
%   \documentclass[tikz,border=4pt]{standalone}
%   % ==========================================================================
% tikz-tcee.tex — Estilos comunes de los gráficos TikZ del temario
% Autor: Víctor Gutiérrez Marcos
%
% Cada gráfico es un fichero standalone en fuentes/<tema>/graficos/:
%
%   \documentclass[tikz,border=4pt]{standalone}
%   % ==========================================================================
% tikz-tcee.tex — Estilos comunes de los gráficos TikZ del temario
% Autor: Víctor Gutiérrez Marcos
%
% Cada gráfico es un fichero standalone en fuentes/<tema>/graficos/:
%
%   \documentclass[tikz,border=4pt]{standalone}
%   \input{../../../latex/tikz-tcee}
%   \begin{document}
%   \begin{tikzpicture}[tcee]
%     \ejes{Q}{P}{6}{5}
%     \draw[curva1] (0.5,4.5) -- (5,0.8) node[etiqueta,right] {$D$};
%     \capa{2}{\draw[curva2] ...;}   % en el vídeo aparece en el paso 2
%   \end{tikzpicture}
%   \end{document}
%
% Paleta: la de la web (morado, dorado, salvia) más tres colores de apoyo
% distinguibles en blanco y negro por el trazo.
% ==========================================================================
\usepackage[T1]{fontenc}
\usepackage{newpxtext,newpxmath}
\usepackage{xcolor}
\usepackage{tikz,pgfplots}
\pgfplotsset{compat=1.18}
\usetikzlibrary{arrows.meta,calc,positioning,patterns,decorations.pathreplacing,intersections,shapes.misc,backgrounds}

\definecolor{tceemorado}{HTML}{5F2987}
\definecolor{tceedorado}{HTML}{B8860B}
\definecolor{tceeverde}{HTML}{3C7D3A}
\definecolor{tceeazul}{HTML}{1F5F99}
\definecolor{tceerojo}{HTML}{B22222}
\definecolor{tceenaranja}{HTML}{E07B28}
\definecolor{tceegris}{HTML}{767171}
\definecolor{tceesalvia}{HTML}{E2EFD9}

% Capas para el vídeo: \capa{n}{…} dibuja su contenido solo a partir del paso n.
% En el PDF, la web y Word se ve todo; generar-video.py compila el gráfico una
% vez por paso con \def\tceepaso{k} para que cada elemento aparezca al narrarlo.
% Vectores con flecha encima, como en el texto (tcee.sty): \vect{x}, \vect{p}
\providecommand{\vect}[1]{\vec{#1}}
\providecommand{\tceepaso}{1000}
\newcommand{\capa}[2]{\ifnum#1>\tceepaso\relax\else#2\fi}

\tikzset{
  tcee/.style={line cap=round,line join=round,font=\small,>={Stealth[length=2.2mm]}},
  eje/.style={->,thick,black},
  curva1/.style={very thick,tceemorado},
  curva2/.style={very thick,tceeazul},
  curva3/.style={very thick,tceeverde},
  curva4/.style={very thick,tceedorado},
  curvanueva/.style={very thick,tceerojo},
  desplazada/.style={very thick,dashed},
  guia/.style={thin,dashed,tceegris},
  punto/.style={circle,fill=black,inner sep=1.3pt},
  etiqueta/.style={font=\small,inner sep=2pt},
  flecha/.style={->,thick,tceerojo},
  area/.style={fill=tceemorado,fill opacity=0.18,draw=none},
  area2/.style={fill=tceedorado,fill opacity=0.22,draw=none},
}

% \ejes{etiqueta x}{etiqueta y}{largo x}{alto y}
\newcommand{\ejes}[4]{%
  \draw[eje] (0,0) -- (#3,0) node[below left=2pt and 0pt] {#1};
  \draw[eje] (0,0) -- (0,#4) node[left] {#2};
  \node[below left] at (0,0) {$0$};}
% \proyeccion{(x,y)}{etq x}{etq y}: líneas guía a los ejes con etiquetas
\newcommand{\proyeccion}[3]{%
  \draw[guia] let \p1=#1 in (\x1,0) node[below] {#2} |- (0,\y1) node[left] {#3};}

% --------------------------------------------------------------------------
% Paleta «dictamen» (dic…): gráficos de datos de la Guía del dictamen (1D01) y
% de cualquier gráfico de coyuntura que quiera los mismos colores. Un concepto,
% un color en todos los gráficos (CRITERIOS §7, «tres reglas», 24/10/2026);
% tabla completa en fuentes/1D01/_trabajo/paleta-1D01.md. No cambia nada de
% lo anterior: son colores y estilos nuevos.
% Actividad: el PIB (y sus variantes: PIB por persona, PIB mundial), morado.
\definecolor{dicpib}{HTML}{5F2987}
% Demanda interna, gama cálida: consumo (naranjas: privado oscuro, público
% claro), inversión (verde) y existencias (gris con trama, aparte).
\definecolor{dicconspriv}{HTML}{C76A1F}
\definecolor{dicconspub}{HTML}{F5C690}
\definecolor{dicfbcf}{HTML}{2A6B3C}
\definecolor{dicexistencias}{HTML}{B4B4B4}
% Demanda externa y comercio, gama fría (azul claro, que se separa del morado)
\definecolor{dicexterna}{HTML}{6AA6DA}
% Zonas: zona del euro (azul marino), España (rojo), resto de países (grises)
\definecolor{dicze}{HTML}{17406D}
\definecolor{dicespana}{HTML}{C8382E}
\definecolor{dicotro}{HTML}{8A8A8A}
\definecolor{dicotroclaro}{HTML}{B5B5B5}
\definecolor{dicotromasclaro}{HTML}{DADADA}
% Precios: inflación general (negro, continuo) y subyacente (negro,
% discontinuo); componentes: energía (bermellón), alimentos (amarillo) y los
% de la subyacente en azules (bienes industriales claro, servicios oscuro).
\definecolor{dicinflacion}{HTML}{1A1A1A}
\definecolor{dicenergia}{HTML}{D55E00}
\definecolor{dicalimentos}{HTML}{F2D45C}
\definecolor{dicbine}{HTML}{7FB2E0}
\definecolor{dicservicios}{HTML}{1F4E89}
% Trabajo: empleo (verde azulado), horas (el mismo, claro), tasa de actividad
% (verde azulado oscuro), paro (teja); productividad (dorado de la web).
\definecolor{dicempleo}{HTML}{1B7F6B}
\definecolor{dichoras}{HTML}{86CBB8}
\definecolor{dicactividad}{HTML}{0F5247}
\definecolor{dicparo}{HTML}{A23B2A}
\definecolor{dicproductividad}{HTML}{B8860B}
% Hogares: renta disponible (pizarra) y tasa de ahorro (negro); el consumo es
% dicconspriv.
\definecolor{dicrenta}{HTML}{4A5A6A}
\definecolor{dicahorro}{HTML}{1A1A1A}
% Política monetaria: tipos del BCE en azules (depósito, el de la zona del
% euro; operaciones principales, medio; facilidad marginal, claro); balance:
% préstamos a los bancos (dorado claro), resto del activo (gris claro).
\definecolor{dicdeposito}{HTML}{17406D}
\definecolor{dicopf}{HTML}{4F86C6}
\definecolor{dicfmc}{HTML}{9EC5EA}
\definecolor{dicprestamos}{HTML}{D9A520}
\definecolor{dicresto}{HTML}{CFCFCF}
% Ciclo: series adelantadas, coincidentes (el PIB) y retrasadas
\definecolor{dicadelantada}{HTML}{D55E00}
\definecolor{dicretrasada}{HTML}{6AA6DA}
% Referencias: rejilla, previsión, sombreados
\definecolor{dicreferencia}{HTML}{767171}

\tikzset{
  % barra de un componente: \addplot[dicbarra=dicfbcf]
  dicbarra/.style={fill=#1, draw=#1!65!black, line width=0.3pt},
  % existencias: gris con trama, para que se distinga también en gris
  dicbarratrama/.style={fill=#1, draw=#1!65!black, line width=0.3pt,
    postaction={pattern=north east lines, pattern color=white}},
  % línea de una serie; con halo blanco cuando pasa sobre barras
  diclinea/.style={very thick, #1},
  diclineahalo/.style={very thick, #1, preaction={draw=white, line width=3.4pt, solid}},
  % puntos de una serie (capa aparte, después de la línea)
  dicpuntos/.style={only marks, mark=*, mark size=1.8pt,
    mark options={fill=white, draw=#1, line width=0.9pt, solid}},
  % línea del cero o de un umbral, después de las barras
  diccero/.style={thin, black},
  dicumbral/.style={thin, dashed, dicreferencia},
}
% Capas de los gráficos de datos (CRITERIOS §7, «tres reglas»), de abajo
% arriba: cuadrícula → sombreados (dicfondo) → barras y áreas (dicbarras) → ejes, cero y umbrales
% (axis lines) → líneas (main) → puntos (dicpuntos) → rótulos (axis
% descriptions). No basta el orden del código: pgfplots dibuja cuadrícula,
% ejes y marcas en sus propias capas (y area style pone los ejes encima), y
% dentro de un axis los \draw, \fill y \node no admiten capa (van con los
% datos de main); por eso todo lo que se dibuja dentro del axis es un \addplot.
% Uso: \begin{axis}[…, dicapas] (al final de las opciones);
%   barras y áreas: \addplot[on layer=dicbarras, …];
%   sombreados:     \dicsombra[estilo]{x1}{y1}{x2}{y2} (rectángulo, en dicfondo,
%                   bajo las barras y sobre la cuadrícula);
%   cero y umbrales: \dicraya[estilo]{(x1,y1) (x2,y2)} (en axis lines);
%   líneas: \addplot normal (main); puntos: \addplot[on layer=dicpuntos, …]
%   (las marcas de cualquier serie van solas a dicpuntos);
%   rótulos: \dicrotulo[opciones del nodo]{(x,y)}{texto} (en axis descriptions),
%   siempre sin fondo (nada de fill=white: si chocan con una línea, se desplazan;
%   CRITERIOS §7, regla 4).
% Los \capa{n} del vídeo no cambian: siguen envolviendo instrucciones completas.
\pgfdeclarelayer{dicfondo}
\pgfdeclarelayer{dicbarras}
\pgfdeclarelayer{dicpuntos}
\pgfplotsset{
  % (el segundo argumento reparte cuadrícula, ejes y rótulos como el estándar)
  layers/dic/.define layer set={axis background, axis grid, dicfondo, dicbarras,
    axis ticks, axis lines, axis tick labels, pre main, main, dicpuntos,
    axis descriptions, axis foreground}{/pgfplots/layers/standard},
  % leyenda sin caja blanca: toma el fondo del gráfico, también si el gráfico va
  % dentro de un cuadro gris (CRITERIOS §7, regla 4, 26/10/2026)
  dicapas/.style={axis on top=false, set layers=dic, mark layer=dicpuntos,
    legend style={fill=none}},
  dicsuelto/.style={forget plot, stack plots=false, sharp plot, mark=none},
}
\newcommand{\dicraya}[2][]{%
  \addplot[on layer=axis lines, dicsuelto, diccero, #1] coordinates {#2};}
\newcommand{\dicsombra}[5][]{%
  \addplot[on layer=dicfondo, dicsuelto, draw=none, #1]
    coordinates {(#2,#3) (#4,#3) (#4,#5) (#2,#5)};}
\newcommand{\dicrotulo}[3][]{%
  \addplot[on layer=axis descriptions, dicsuelto, draw=none] coordinates {#2}
    node[pos=0, #1] {#3};}

%   \begin{document}
%   \begin{tikzpicture}[tcee]
%     \ejes{Q}{P}{6}{5}
%     \draw[curva1] (0.5,4.5) -- (5,0.8) node[etiqueta,right] {$D$};
%     \capa{2}{\draw[curva2] ...;}   % en el vídeo aparece en el paso 2
%   \end{tikzpicture}
%   \end{document}
%
% Paleta: la de la web (morado, dorado, salvia) más tres colores de apoyo
% distinguibles en blanco y negro por el trazo.
% ==========================================================================
\usepackage[T1]{fontenc}
\usepackage{newpxtext,newpxmath}
\usepackage{xcolor}
\usepackage{tikz,pgfplots}
\pgfplotsset{compat=1.18}
\usetikzlibrary{arrows.meta,calc,positioning,patterns,decorations.pathreplacing,intersections,shapes.misc,backgrounds}

\definecolor{tceemorado}{HTML}{5F2987}
\definecolor{tceedorado}{HTML}{B8860B}
\definecolor{tceeverde}{HTML}{3C7D3A}
\definecolor{tceeazul}{HTML}{1F5F99}
\definecolor{tceerojo}{HTML}{B22222}
\definecolor{tceenaranja}{HTML}{E07B28}
\definecolor{tceegris}{HTML}{767171}
\definecolor{tceesalvia}{HTML}{E2EFD9}

% Capas para el vídeo: \capa{n}{…} dibuja su contenido solo a partir del paso n.
% En el PDF, la web y Word se ve todo; generar-video.py compila el gráfico una
% vez por paso con \def\tceepaso{k} para que cada elemento aparezca al narrarlo.
% Vectores con flecha encima, como en el texto (tcee.sty): \vect{x}, \vect{p}
\providecommand{\vect}[1]{\vec{#1}}
\providecommand{\tceepaso}{1000}
\newcommand{\capa}[2]{\ifnum#1>\tceepaso\relax\else#2\fi}

\tikzset{
  tcee/.style={line cap=round,line join=round,font=\small,>={Stealth[length=2.2mm]}},
  eje/.style={->,thick,black},
  curva1/.style={very thick,tceemorado},
  curva2/.style={very thick,tceeazul},
  curva3/.style={very thick,tceeverde},
  curva4/.style={very thick,tceedorado},
  curvanueva/.style={very thick,tceerojo},
  desplazada/.style={very thick,dashed},
  guia/.style={thin,dashed,tceegris},
  punto/.style={circle,fill=black,inner sep=1.3pt},
  etiqueta/.style={font=\small,inner sep=2pt},
  flecha/.style={->,thick,tceerojo},
  area/.style={fill=tceemorado,fill opacity=0.18,draw=none},
  area2/.style={fill=tceedorado,fill opacity=0.22,draw=none},
}

% \ejes{etiqueta x}{etiqueta y}{largo x}{alto y}
\newcommand{\ejes}[4]{%
  \draw[eje] (0,0) -- (#3,0) node[below left=2pt and 0pt] {#1};
  \draw[eje] (0,0) -- (0,#4) node[left] {#2};
  \node[below left] at (0,0) {$0$};}
% \proyeccion{(x,y)}{etq x}{etq y}: líneas guía a los ejes con etiquetas
\newcommand{\proyeccion}[3]{%
  \draw[guia] let \p1=#1 in (\x1,0) node[below] {#2} |- (0,\y1) node[left] {#3};}

% --------------------------------------------------------------------------
% Paleta «dictamen» (dic…): gráficos de datos de la Guía del dictamen (1D01) y
% de cualquier gráfico de coyuntura que quiera los mismos colores. Un concepto,
% un color en todos los gráficos (CRITERIOS §7, «tres reglas», 24/10/2026);
% tabla completa en fuentes/1D01/_trabajo/paleta-1D01.md. No cambia nada de
% lo anterior: son colores y estilos nuevos.
% Actividad: el PIB (y sus variantes: PIB por persona, PIB mundial), morado.
\definecolor{dicpib}{HTML}{5F2987}
% Demanda interna, gama cálida: consumo (naranjas: privado oscuro, público
% claro), inversión (verde) y existencias (gris con trama, aparte).
\definecolor{dicconspriv}{HTML}{C76A1F}
\definecolor{dicconspub}{HTML}{F5C690}
\definecolor{dicfbcf}{HTML}{2A6B3C}
\definecolor{dicexistencias}{HTML}{B4B4B4}
% Demanda externa y comercio, gama fría (azul claro, que se separa del morado)
\definecolor{dicexterna}{HTML}{6AA6DA}
% Zonas: zona del euro (azul marino), España (rojo), resto de países (grises)
\definecolor{dicze}{HTML}{17406D}
\definecolor{dicespana}{HTML}{C8382E}
\definecolor{dicotro}{HTML}{8A8A8A}
\definecolor{dicotroclaro}{HTML}{B5B5B5}
\definecolor{dicotromasclaro}{HTML}{DADADA}
% Precios: inflación general (negro, continuo) y subyacente (negro,
% discontinuo); componentes: energía (bermellón), alimentos (amarillo) y los
% de la subyacente en azules (bienes industriales claro, servicios oscuro).
\definecolor{dicinflacion}{HTML}{1A1A1A}
\definecolor{dicenergia}{HTML}{D55E00}
\definecolor{dicalimentos}{HTML}{F2D45C}
\definecolor{dicbine}{HTML}{7FB2E0}
\definecolor{dicservicios}{HTML}{1F4E89}
% Trabajo: empleo (verde azulado), horas (el mismo, claro), tasa de actividad
% (verde azulado oscuro), paro (teja); productividad (dorado de la web).
\definecolor{dicempleo}{HTML}{1B7F6B}
\definecolor{dichoras}{HTML}{86CBB8}
\definecolor{dicactividad}{HTML}{0F5247}
\definecolor{dicparo}{HTML}{A23B2A}
\definecolor{dicproductividad}{HTML}{B8860B}
% Hogares: renta disponible (pizarra) y tasa de ahorro (negro); el consumo es
% dicconspriv.
\definecolor{dicrenta}{HTML}{4A5A6A}
\definecolor{dicahorro}{HTML}{1A1A1A}
% Política monetaria: tipos del BCE en azules (depósito, el de la zona del
% euro; operaciones principales, medio; facilidad marginal, claro); balance:
% préstamos a los bancos (dorado claro), resto del activo (gris claro).
\definecolor{dicdeposito}{HTML}{17406D}
\definecolor{dicopf}{HTML}{4F86C6}
\definecolor{dicfmc}{HTML}{9EC5EA}
\definecolor{dicprestamos}{HTML}{D9A520}
\definecolor{dicresto}{HTML}{CFCFCF}
% Ciclo: series adelantadas, coincidentes (el PIB) y retrasadas
\definecolor{dicadelantada}{HTML}{D55E00}
\definecolor{dicretrasada}{HTML}{6AA6DA}
% Referencias: rejilla, previsión, sombreados
\definecolor{dicreferencia}{HTML}{767171}

\tikzset{
  % barra de un componente: \addplot[dicbarra=dicfbcf]
  dicbarra/.style={fill=#1, draw=#1!65!black, line width=0.3pt},
  % existencias: gris con trama, para que se distinga también en gris
  dicbarratrama/.style={fill=#1, draw=#1!65!black, line width=0.3pt,
    postaction={pattern=north east lines, pattern color=white}},
  % línea de una serie; con halo blanco cuando pasa sobre barras
  diclinea/.style={very thick, #1},
  diclineahalo/.style={very thick, #1, preaction={draw=white, line width=3.4pt, solid}},
  % puntos de una serie (capa aparte, después de la línea)
  dicpuntos/.style={only marks, mark=*, mark size=1.8pt,
    mark options={fill=white, draw=#1, line width=0.9pt, solid}},
  % línea del cero o de un umbral, después de las barras
  diccero/.style={thin, black},
  dicumbral/.style={thin, dashed, dicreferencia},
}
% Capas de los gráficos de datos (CRITERIOS §7, «tres reglas»), de abajo
% arriba: cuadrícula → sombreados (dicfondo) → barras y áreas (dicbarras) → ejes, cero y umbrales
% (axis lines) → líneas (main) → puntos (dicpuntos) → rótulos (axis
% descriptions). No basta el orden del código: pgfplots dibuja cuadrícula,
% ejes y marcas en sus propias capas (y area style pone los ejes encima), y
% dentro de un axis los \draw, \fill y \node no admiten capa (van con los
% datos de main); por eso todo lo que se dibuja dentro del axis es un \addplot.
% Uso: \begin{axis}[…, dicapas] (al final de las opciones);
%   barras y áreas: \addplot[on layer=dicbarras, …];
%   sombreados:     \dicsombra[estilo]{x1}{y1}{x2}{y2} (rectángulo, en dicfondo,
%                   bajo las barras y sobre la cuadrícula);
%   cero y umbrales: \dicraya[estilo]{(x1,y1) (x2,y2)} (en axis lines);
%   líneas: \addplot normal (main); puntos: \addplot[on layer=dicpuntos, …]
%   (las marcas de cualquier serie van solas a dicpuntos);
%   rótulos: \dicrotulo[opciones del nodo]{(x,y)}{texto} (en axis descriptions),
%   siempre sin fondo (nada de fill=white: si chocan con una línea, se desplazan;
%   CRITERIOS §7, regla 4).
% Los \capa{n} del vídeo no cambian: siguen envolviendo instrucciones completas.
\pgfdeclarelayer{dicfondo}
\pgfdeclarelayer{dicbarras}
\pgfdeclarelayer{dicpuntos}
\pgfplotsset{
  % (el segundo argumento reparte cuadrícula, ejes y rótulos como el estándar)
  layers/dic/.define layer set={axis background, axis grid, dicfondo, dicbarras,
    axis ticks, axis lines, axis tick labels, pre main, main, dicpuntos,
    axis descriptions, axis foreground}{/pgfplots/layers/standard},
  % leyenda sin caja blanca: toma el fondo del gráfico, también si el gráfico va
  % dentro de un cuadro gris (CRITERIOS §7, regla 4, 26/10/2026)
  dicapas/.style={axis on top=false, set layers=dic, mark layer=dicpuntos,
    legend style={fill=none}},
  dicsuelto/.style={forget plot, stack plots=false, sharp plot, mark=none},
}
\newcommand{\dicraya}[2][]{%
  \addplot[on layer=axis lines, dicsuelto, diccero, #1] coordinates {#2};}
\newcommand{\dicsombra}[5][]{%
  \addplot[on layer=dicfondo, dicsuelto, draw=none, #1]
    coordinates {(#2,#3) (#4,#3) (#4,#5) (#2,#5)};}
\newcommand{\dicrotulo}[3][]{%
  \addplot[on layer=axis descriptions, dicsuelto, draw=none] coordinates {#2}
    node[pos=0, #1] {#3};}

%   \begin{document}
%   \begin{tikzpicture}[tcee]
%     \ejes{Q}{P}{6}{5}
%     \draw[curva1] (0.5,4.5) -- (5,0.8) node[etiqueta,right] {$D$};
%     \capa{2}{\draw[curva2] ...;}   % en el vídeo aparece en el paso 2
%   \end{tikzpicture}
%   \end{document}
%
% Paleta: la de la web (morado, dorado, salvia) más tres colores de apoyo
% distinguibles en blanco y negro por el trazo.
% ==========================================================================
\usepackage[T1]{fontenc}
\usepackage{newpxtext,newpxmath}
\usepackage{xcolor}
\usepackage{tikz,pgfplots}
\pgfplotsset{compat=1.18}
\usetikzlibrary{arrows.meta,calc,positioning,patterns,decorations.pathreplacing,intersections,shapes.misc,backgrounds}

\definecolor{tceemorado}{HTML}{5F2987}
\definecolor{tceedorado}{HTML}{B8860B}
\definecolor{tceeverde}{HTML}{3C7D3A}
\definecolor{tceeazul}{HTML}{1F5F99}
\definecolor{tceerojo}{HTML}{B22222}
\definecolor{tceenaranja}{HTML}{E07B28}
\definecolor{tceegris}{HTML}{767171}
\definecolor{tceesalvia}{HTML}{E2EFD9}

% Capas para el vídeo: \capa{n}{…} dibuja su contenido solo a partir del paso n.
% En el PDF, la web y Word se ve todo; generar-video.py compila el gráfico una
% vez por paso con \def\tceepaso{k} para que cada elemento aparezca al narrarlo.
% Vectores con flecha encima, como en el texto (tcee.sty): \vect{x}, \vect{p}
\providecommand{\vect}[1]{\vec{#1}}
\providecommand{\tceepaso}{1000}
\newcommand{\capa}[2]{\ifnum#1>\tceepaso\relax\else#2\fi}

\tikzset{
  tcee/.style={line cap=round,line join=round,font=\small,>={Stealth[length=2.2mm]}},
  eje/.style={->,thick,black},
  curva1/.style={very thick,tceemorado},
  curva2/.style={very thick,tceeazul},
  curva3/.style={very thick,tceeverde},
  curva4/.style={very thick,tceedorado},
  curvanueva/.style={very thick,tceerojo},
  desplazada/.style={very thick,dashed},
  guia/.style={thin,dashed,tceegris},
  punto/.style={circle,fill=black,inner sep=1.3pt},
  etiqueta/.style={font=\small,inner sep=2pt},
  flecha/.style={->,thick,tceerojo},
  area/.style={fill=tceemorado,fill opacity=0.18,draw=none},
  area2/.style={fill=tceedorado,fill opacity=0.22,draw=none},
}

% \ejes{etiqueta x}{etiqueta y}{largo x}{alto y}
\newcommand{\ejes}[4]{%
  \draw[eje] (0,0) -- (#3,0) node[below left=2pt and 0pt] {#1};
  \draw[eje] (0,0) -- (0,#4) node[left] {#2};
  \node[below left] at (0,0) {$0$};}
% \proyeccion{(x,y)}{etq x}{etq y}: líneas guía a los ejes con etiquetas
\newcommand{\proyeccion}[3]{%
  \draw[guia] let \p1=#1 in (\x1,0) node[below] {#2} |- (0,\y1) node[left] {#3};}

% --------------------------------------------------------------------------
% Paleta «dictamen» (dic…): gráficos de datos de la Guía del dictamen (1D01) y
% de cualquier gráfico de coyuntura que quiera los mismos colores. Un concepto,
% un color en todos los gráficos (CRITERIOS §7, «tres reglas», 24/10/2026);
% tabla completa en fuentes/1D01/_trabajo/paleta-1D01.md. No cambia nada de
% lo anterior: son colores y estilos nuevos.
% Actividad: el PIB (y sus variantes: PIB por persona, PIB mundial), morado.
\definecolor{dicpib}{HTML}{5F2987}
% Demanda interna, gama cálida: consumo (naranjas: privado oscuro, público
% claro), inversión (verde) y existencias (gris con trama, aparte).
\definecolor{dicconspriv}{HTML}{C76A1F}
\definecolor{dicconspub}{HTML}{F5C690}
\definecolor{dicfbcf}{HTML}{2A6B3C}
\definecolor{dicexistencias}{HTML}{B4B4B4}
% Demanda externa y comercio, gama fría (azul claro, que se separa del morado)
\definecolor{dicexterna}{HTML}{6AA6DA}
% Zonas: zona del euro (azul marino), España (rojo), resto de países (grises)
\definecolor{dicze}{HTML}{17406D}
\definecolor{dicespana}{HTML}{C8382E}
\definecolor{dicotro}{HTML}{8A8A8A}
\definecolor{dicotroclaro}{HTML}{B5B5B5}
\definecolor{dicotromasclaro}{HTML}{DADADA}
% Precios: inflación general (negro, continuo) y subyacente (negro,
% discontinuo); componentes: energía (bermellón), alimentos (amarillo) y los
% de la subyacente en azules (bienes industriales claro, servicios oscuro).
\definecolor{dicinflacion}{HTML}{1A1A1A}
\definecolor{dicenergia}{HTML}{D55E00}
\definecolor{dicalimentos}{HTML}{F2D45C}
\definecolor{dicbine}{HTML}{7FB2E0}
\definecolor{dicservicios}{HTML}{1F4E89}
% Trabajo: empleo (verde azulado), horas (el mismo, claro), tasa de actividad
% (verde azulado oscuro), paro (teja); productividad (dorado de la web).
\definecolor{dicempleo}{HTML}{1B7F6B}
\definecolor{dichoras}{HTML}{86CBB8}
\definecolor{dicactividad}{HTML}{0F5247}
\definecolor{dicparo}{HTML}{A23B2A}
\definecolor{dicproductividad}{HTML}{B8860B}
% Hogares: renta disponible (pizarra) y tasa de ahorro (negro); el consumo es
% dicconspriv.
\definecolor{dicrenta}{HTML}{4A5A6A}
\definecolor{dicahorro}{HTML}{1A1A1A}
% Política monetaria: tipos del BCE en azules (depósito, el de la zona del
% euro; operaciones principales, medio; facilidad marginal, claro); balance:
% préstamos a los bancos (dorado claro), resto del activo (gris claro).
\definecolor{dicdeposito}{HTML}{17406D}
\definecolor{dicopf}{HTML}{4F86C6}
\definecolor{dicfmc}{HTML}{9EC5EA}
\definecolor{dicprestamos}{HTML}{D9A520}
\definecolor{dicresto}{HTML}{CFCFCF}
% Ciclo: series adelantadas, coincidentes (el PIB) y retrasadas
\definecolor{dicadelantada}{HTML}{D55E00}
\definecolor{dicretrasada}{HTML}{6AA6DA}
% Referencias: rejilla, previsión, sombreados
\definecolor{dicreferencia}{HTML}{767171}

\tikzset{
  % barra de un componente: \addplot[dicbarra=dicfbcf]
  dicbarra/.style={fill=#1, draw=#1!65!black, line width=0.3pt},
  % existencias: gris con trama, para que se distinga también en gris
  dicbarratrama/.style={fill=#1, draw=#1!65!black, line width=0.3pt,
    postaction={pattern=north east lines, pattern color=white}},
  % línea de una serie; con halo blanco cuando pasa sobre barras
  diclinea/.style={very thick, #1},
  diclineahalo/.style={very thick, #1, preaction={draw=white, line width=3.4pt, solid}},
  % puntos de una serie (capa aparte, después de la línea)
  dicpuntos/.style={only marks, mark=*, mark size=1.8pt,
    mark options={fill=white, draw=#1, line width=0.9pt, solid}},
  % línea del cero o de un umbral, después de las barras
  diccero/.style={thin, black},
  dicumbral/.style={thin, dashed, dicreferencia},
}
% Capas de los gráficos de datos (CRITERIOS §7, «tres reglas»), de abajo
% arriba: cuadrícula → sombreados (dicfondo) → barras y áreas (dicbarras) → ejes, cero y umbrales
% (axis lines) → líneas (main) → puntos (dicpuntos) → rótulos (axis
% descriptions). No basta el orden del código: pgfplots dibuja cuadrícula,
% ejes y marcas en sus propias capas (y area style pone los ejes encima), y
% dentro de un axis los \draw, \fill y \node no admiten capa (van con los
% datos de main); por eso todo lo que se dibuja dentro del axis es un \addplot.
% Uso: \begin{axis}[…, dicapas] (al final de las opciones);
%   barras y áreas: \addplot[on layer=dicbarras, …];
%   sombreados:     \dicsombra[estilo]{x1}{y1}{x2}{y2} (rectángulo, en dicfondo,
%                   bajo las barras y sobre la cuadrícula);
%   cero y umbrales: \dicraya[estilo]{(x1,y1) (x2,y2)} (en axis lines);
%   líneas: \addplot normal (main); puntos: \addplot[on layer=dicpuntos, …]
%   (las marcas de cualquier serie van solas a dicpuntos);
%   rótulos: \dicrotulo[opciones del nodo]{(x,y)}{texto} (en axis descriptions),
%   siempre sin fondo (nada de fill=white: si chocan con una línea, se desplazan;
%   CRITERIOS §7, regla 4).
% Los \capa{n} del vídeo no cambian: siguen envolviendo instrucciones completas.
\pgfdeclarelayer{dicfondo}
\pgfdeclarelayer{dicbarras}
\pgfdeclarelayer{dicpuntos}
\pgfplotsset{
  % (el segundo argumento reparte cuadrícula, ejes y rótulos como el estándar)
  layers/dic/.define layer set={axis background, axis grid, dicfondo, dicbarras,
    axis ticks, axis lines, axis tick labels, pre main, main, dicpuntos,
    axis descriptions, axis foreground}{/pgfplots/layers/standard},
  % leyenda sin caja blanca: toma el fondo del gráfico, también si el gráfico va
  % dentro de un cuadro gris (CRITERIOS §7, regla 4, 26/10/2026)
  dicapas/.style={axis on top=false, set layers=dic, mark layer=dicpuntos,
    legend style={fill=none}},
  dicsuelto/.style={forget plot, stack plots=false, sharp plot, mark=none},
}
\newcommand{\dicraya}[2][]{%
  \addplot[on layer=axis lines, dicsuelto, diccero, #1] coordinates {#2};}
\newcommand{\dicsombra}[5][]{%
  \addplot[on layer=dicfondo, dicsuelto, draw=none, #1]
    coordinates {(#2,#3) (#4,#3) (#4,#5) (#2,#5)};}
\newcommand{\dicrotulo}[3][]{%
  \addplot[on layer=axis descriptions, dicsuelto, draw=none] coordinates {#2}
    node[pos=0, #1] {#3};}

\begin{document}
\begin{tikzpicture}[tcee,
    declare function={
      fpp(\x)=6*sqrt(1-(\x/7)^2);
      cc(\x)=1.5*\x*4.199125/((2/3)*(5-\x)+1.5*\x);
      ua(\x,\u)=(\u/\x^0.4)^(1/0.6);
      ub(\x,\v)=4.199125-(\v/(5-\x)^0.6)^(1/0.4);
      fpuu(\t,\k)=7*cos(\k)^(0.2)*cos(\t)^(0.8);
      fpuv(\t,\k)=7*sin(\k)^(0.2)*sin(\t)^(0.8);}]
  % 1: ejes X, Y
  \capa{1}{\ejes{$X$}{$Y$}{8.4}{7}}
  % 2: la FPP
  \capa{2}{\draw[curva2] plot[domain=90:0,samples=150,variable=\t] ({7*cos(\t)},{6*sin(\t)});}
  \capa{2}{\node[etiqueta,text=tceeazul,anchor=south west] at (6.75,1.75) {FPP};}
  % 3: punto E y sus guías X*, Y*
  \capa{3}{\node[punto] (E) at (5,4.199125) {};}
  \capa{3}{\node[etiqueta,anchor=south west,inner xsep=0.5pt] (Elab) at (E.north east) {$E$};}
  \capa{3}{\proyeccion{(E)}{$X^*$}{$Y^*$}}
  % 4: la caja del consumo, con O_A en el origen y O_B en E
  \capa{4}{\draw[thick] (0,0) rectangle (5,4.199125);}
  \capa{4}{\node[etiqueta,text=tceemorado,anchor=south west] at (0.03,0.03) {$O_A$};}
  \capa{4}{\node[etiqueta,text=tceeverde,anchor=north east] at (4.97,4.17) {$O_B$};}
  % 5: curva de contrato dentro de la caja
  \capa{5}{\draw[curva2,thick] plot[domain=0:5,samples=120] (\x,{cc(\x)});}
  \capa{5}{\node[etiqueta,text=tceeazul,anchor=north west,align=left] at (3.3,3.1) {curva de\\ contrato};}
  % 6: recta tangente a la FPP en E (RMT)
  \capa{6}{\draw[curva4] (3.3,5.6863) -- (7.9,1.6623) node[etiqueta,right] {$\mathrm{RMT}_{XY}$};}
  % 7: E*, curvas de indiferencia tangentes y la paralela (RMS = RMT)
  \capa{7}{\draw[curva1,thick] plot[domain=1.15:3.0,samples=60] (\x,{ua(\x,2.0716576)});}
  \capa{7}{\draw[curva3,thick] plot[domain=0.6:3.0,samples=60] (\x,{ub(\x,2.6114129)});}
  \capa{7}{\draw[curva4,desplazada,thick] (0.8,3.1493) -- (3.05,1.1810);}
  \capa{7}{\node[punto,label={[etiqueta]above right:$E^*$}] at (1.76,2.309519) {};}
  \capa{7}{\node[etiqueta,text=tceedorado,anchor=north west,align=left,font=\footnotesize] at (0.12,4.12) {$\mathrm{RMS}^A_{XY}=\mathrm{RMS}^B_{XY}$\\ ${}=\mathrm{RMT}_{XY}$};}
  % Panel derecho: el espacio de las utilidades
  \begin{scope}[shift={(10.2,0)},x=0.85cm,y=0.85cm]
    % 8: ejes U_A, U_B
    \capa{8}{\ejes{$U_A$}{$U_B$}{8.4}{8.2}}
    % 9: FPU_1 con E*_1 (y E_1 en la FPP)
    \capa{9}{\draw[curva1,thick] plot[smooth,tension=0.3] coordinates {(0.000,7.500) (0.079,7.500) (0.157,7.500) (0.236,7.500) (0.314,7.500) (0.393,7.500) (0.472,7.500) (0.551,7.500) (0.630,7.499) (0.709,7.499) (0.788,7.498) (0.867,7.497) (0.947,7.496) (1.027,7.494) (1.106,7.492) (1.186,7.489) (1.266,7.486) (1.346,7.482) (1.426,7.476) (1.506,7.470) (1.586,7.463) (1.666,7.454) (1.746,7.444) (1.826,7.433) (1.905,7.419) (1.984,7.404) (2.062,7.387) (2.140,7.367) (2.218,7.345) (2.294,7.320) (2.370,7.293) (2.444,7.263) (2.518,7.229) (2.590,7.193) (2.660,7.153) (2.729,7.110) (2.797,7.064) (2.862,7.014) (2.926,6.961) (2.988,6.904) (3.047,6.844) (3.105,6.781) (3.160,6.715) (3.213,6.646) (3.263,6.574) (3.311,6.499) (3.357,6.421) (3.400,6.342) (3.442,6.260) (3.480,6.177) (3.517,6.092) (3.551,6.005) (3.584,5.918) (3.614,5.829) (3.642,5.740) (3.669,5.650) (3.694,5.559) (3.717,5.469) (3.738,5.379) (3.758,5.288) (3.777,5.198) (3.794,5.108) (3.810,5.019) (3.825,4.931) (3.838,4.843) (3.851,4.755) (3.863,4.669) (3.873,4.583) (3.883,4.499) (3.893,4.415) (3.901,4.333) (3.909,4.251) (3.916,4.170) (3.923,4.091) (3.929,4.012) (3.935,3.935) (3.940,3.859) (3.945,3.783) (3.950,3.709) (3.954,3.636) (3.957,3.563) (3.961,3.492) (3.964,3.422) (3.967,3.353) (3.970,3.284) (3.972,3.217) (3.975,3.150) (3.977,3.085) (3.979,3.020) (3.981,2.956) (3.982,2.893) (3.984,2.831) (3.985,2.770) (3.987,2.709) (3.988,2.650) (3.989,2.590) (3.990,2.532) (3.991,2.474) (3.992,2.418) (3.993,2.361) (3.993,2.306) (3.994,2.250) (3.995,2.196) (3.995,2.142) (3.996,2.089) (3.996,2.036) (3.996,1.984) (3.997,1.932) (3.997,1.881) (3.997,1.830) (3.998,1.780) (3.998,1.730) (3.998,1.681) (3.998,1.632) (3.999,1.583) (3.999,1.535) (3.999,1.487) (3.999,1.440) (3.999,1.393) (3.999,1.346) (3.999,1.299) (3.999,1.253) (4.000,1.208) (4.000,1.162) (4.000,1.117) (4.000,1.072) (4.000,1.027) (4.000,0.982) (4.000,0.938) (4.000,0.894) (4.000,0.850) (4.000,0.807) (4.000,0.763) (4.000,0.720) (4.000,0.677) (4.000,0.634) (4.000,0.591) (4.000,0.548) (4.000,0.505) (4.000,0.463) (4.000,0.420) (4.000,0.378) (4.000,0.336) (4.000,0.294) (4.000,0.252) (4.000,0.210) (4.000,0.168) (4.000,0.126) (4.000,0.084) (4.000,0.042) (4.000,0.000)}
      node[etiqueta,below] {$\mathrm{FPU}_1$};}
    \capa{9}{\node[punto,label={[etiqueta]above right:$E^*_1$}] at (2.7294,7.1102) {};}
    % 10: FPU_2, la de la caja de E, con E*_2
    \capa{10}{\draw[curva3,thick] plot[smooth,tension=0.3] coordinates {(0.000,5.900) (0.062,5.900) (0.124,5.900) (0.185,5.900) (0.247,5.900) (0.309,5.900) (0.371,5.900) (0.433,5.900) (0.495,5.900) (0.558,5.900) (0.620,5.900) (0.683,5.900) (0.745,5.900) (0.808,5.900) (0.871,5.900) (0.934,5.900) (0.998,5.899) (1.061,5.899) (1.125,5.899) (1.189,5.899) (1.254,5.898) (1.318,5.898) (1.383,5.897) (1.448,5.897) (1.514,5.896) (1.580,5.895) (1.646,5.895) (1.712,5.894) (1.779,5.893) (1.846,5.891) (1.914,5.890) (1.982,5.888) (2.050,5.886) (2.119,5.884) (2.188,5.882) (2.257,5.880) (2.327,5.877) (2.397,5.874) (2.468,5.870) (2.539,5.867) (2.610,5.863) (2.682,5.858) (2.754,5.853) (2.827,5.847) (2.900,5.841) (2.973,5.835) (3.047,5.828) (3.120,5.820) (3.195,5.811) (3.269,5.802) (3.344,5.792) (3.419,5.781) (3.494,5.769) (3.569,5.756) (3.644,5.743) (3.720,5.728) (3.795,5.712) (3.870,5.695) (3.945,5.676) (4.020,5.657) (4.095,5.636) (4.169,5.613) (4.243,5.589) (4.316,5.564) (4.388,5.537) (4.460,5.508) (4.531,5.478) (4.602,5.445) (4.671,5.411) (4.739,5.376) (4.806,5.338) (4.872,5.298) (4.937,5.257) (5.000,5.214) (5.061,5.169) (5.121,5.121) (5.180,5.072) (5.236,5.022) (5.291,4.969) (5.344,4.914) (5.396,4.858) (5.445,4.800) (5.492,4.741) (5.538,4.680) (5.581,4.617) (5.623,4.553) (5.662,4.488) (5.700,4.421) (5.736,4.353) (5.769,4.285) (5.801,4.215) (5.832,4.144) (5.860,4.073) (5.887,4.001) (5.912,3.928) (5.936,3.855) (5.958,3.781) (5.978,3.707) (5.997,3.632) (6.015,3.557) (6.032,3.482) (6.047,3.407) (6.061,3.332) (6.074,3.257) (6.086,3.182) (6.097,3.107) (6.108,3.032) (6.117,2.957) (6.126,2.883) (6.134,2.808) (6.141,2.734) (6.147,2.660) (6.153,2.587) (6.159,2.513) (6.163,2.440) (6.168,2.368) (6.172,2.295) (6.175,2.223) (6.179,2.152) (6.181,2.080) (6.184,2.009) (6.186,1.939) (6.188,1.868) (6.190,1.798) (6.191,1.729) (6.193,1.659) (6.194,1.590) (6.195,1.522) (6.196,1.453) (6.197,1.385) (6.197,1.317) (6.198,1.250) (6.198,1.182) (6.199,1.115) (6.199,1.048) (6.199,0.982) (6.199,0.915) (6.200,0.849) (6.200,0.783) (6.200,0.717) (6.200,0.652) (6.200,0.586) (6.200,0.521) (6.200,0.455) (6.200,0.390) (6.200,0.325) (6.200,0.260) (6.200,0.195) (6.200,0.130) (6.200,0.065) (6.200,0.000)};}
    \capa{10}{\node[etiqueta,text=tceeverde,left] at (0,5.9000) {$\mathrm{FPU}_2$};}
    \capa{10}{\node[punto,label={[etiqueta]above right:$E^*_2$}] at (5.1214,5.1214) {};}
    % 11: FPU_3 con E*_3
    \capa{11}{\draw[curva4,thick] plot[smooth,tension=0.3] coordinates {(0.000,3.900) (0.041,3.900) (0.082,3.900) (0.123,3.900) (0.163,3.900) (0.204,3.900) (0.245,3.900) (0.286,3.900) (0.327,3.900) (0.369,3.900) (0.410,3.900) (0.451,3.900) (0.493,3.900) (0.534,3.900) (0.576,3.900) (0.618,3.900) (0.660,3.900) (0.702,3.900) (0.744,3.900) (0.786,3.900) (0.829,3.900) (0.872,3.899) (0.915,3.899) (0.958,3.899) (1.001,3.899) (1.045,3.899) (1.089,3.899) (1.133,3.899) (1.177,3.898) (1.222,3.898) (1.267,3.898) (1.312,3.898) (1.357,3.897) (1.403,3.897) (1.449,3.897) (1.496,3.896) (1.542,3.896) (1.590,3.895) (1.637,3.895) (1.685,3.894) (1.734,3.894) (1.782,3.893) (1.832,3.892) (1.881,3.892) (1.931,3.891) (1.982,3.890) (2.033,3.889) (2.085,3.888) (2.137,3.887) (2.189,3.886) (2.243,3.884) (2.296,3.883) (2.351,3.882) (2.406,3.880) (2.461,3.878) (2.517,3.877) (2.574,3.875) (2.632,3.873) (2.690,3.870) (2.749,3.868) (2.808,3.865) (2.869,3.863) (2.930,3.860) (2.992,3.857) (3.054,3.853) (3.118,3.850) (3.182,3.846) (3.247,3.842) (3.313,3.838) (3.379,3.833) (3.447,3.828) (3.516,3.823) (3.585,3.818) (3.655,3.812) (3.726,3.805) (3.798,3.798) (3.871,3.791) (3.945,3.783) (4.020,3.775) (4.096,3.766) (4.172,3.757) (4.250,3.747) (4.328,3.736) (4.408,3.725) (4.488,3.713) (4.569,3.700) (4.651,3.686) (4.733,3.671) (4.816,3.656) (4.900,3.639) (4.984,3.621) (5.069,3.603) (5.155,3.583) (5.241,3.561) (5.327,3.539) (5.413,3.515) (5.499,3.490) (5.585,3.463) (5.671,3.435) (5.757,3.405) (5.842,3.373) (5.927,3.340) (6.011,3.305) (6.094,3.268) (6.176,3.229) (6.257,3.188) (6.336,3.145) (6.413,3.100) (6.489,3.054) (6.563,3.005) (6.635,2.954) (6.705,2.901) (6.772,2.847) (6.837,2.790) (6.899,2.731) (6.958,2.671) (7.014,2.609) (7.068,2.545) (7.118,2.479) (7.166,2.412) (7.211,2.343) (7.253,2.273) (7.291,2.201) (7.327,2.129) (7.360,2.055) (7.390,1.980) (7.418,1.905) (7.443,1.828) (7.465,1.751) (7.486,1.673) (7.504,1.595) (7.519,1.516) (7.533,1.437) (7.545,1.358) (7.556,1.278) (7.565,1.198) (7.572,1.118) (7.579,1.038) (7.584,0.958) (7.588,0.878) (7.592,0.798) (7.594,0.718) (7.596,0.638) (7.598,0.558) (7.599,0.478) (7.599,0.398) (7.600,0.319) (7.600,0.239) (7.600,0.159) (7.600,0.080) (7.600,0.000)};}
    \capa{11}{\node[etiqueta,text=tceedorado,left] at (0,3.9000) {$\mathrm{FPU}_3$};}
    \capa{11}{\node[punto,label={[etiqueta]above right:$E^*_3$}] at (7.1347,2.4567) {};}
    % 12: la GFPU, envolvente de los puntos globalmente eficientes (irregular)
    \capa{12}{\draw[curva2,line width=1.8pt] plot[smooth,tension=0.3] coordinates {(0.000,7.747) (0.081,7.746) (0.162,7.745) (0.243,7.744) (0.324,7.742) (0.406,7.740) (0.487,7.737) (0.568,7.734) (0.649,7.730) (0.730,7.726) (0.811,7.720) (0.893,7.714) (0.974,7.707) (1.055,7.698) (1.135,7.689) (1.216,7.678) (1.296,7.665) (1.377,7.651) (1.456,7.635) (1.536,7.618) (1.615,7.599) (1.694,7.578) (1.772,7.555) (1.850,7.531) (1.927,7.505) (2.003,7.477) (2.079,7.448) (2.155,7.416) (2.229,7.384) (2.303,7.349) (2.376,7.313) (2.449,7.276) (2.520,7.237) (2.591,7.196) (2.661,7.154) (2.729,7.110) (2.797,7.065) (2.864,7.017) (2.929,6.968) (2.993,6.917) (3.056,6.864) (3.117,6.808) (3.177,6.751) (3.234,6.691) (3.290,6.628) (3.344,6.563) (3.396,6.496) (3.445,6.426) (3.493,6.353) (3.537,6.278) (3.582,6.204) (3.629,6.136) (3.677,6.072) (3.728,6.012) (3.781,5.958) (3.838,5.909) (3.897,5.865) (3.960,5.827) (4.026,5.793) (4.096,5.764) (4.168,5.736) (4.238,5.706) (4.306,5.674) (4.374,5.639) (4.441,5.603) (4.506,5.564) (4.571,5.525) (4.634,5.484) (4.697,5.441) (4.759,5.398) (4.821,5.354) (4.882,5.309) (4.942,5.263) (5.002,5.216) (5.062,5.169) (5.121,5.121) (5.180,5.073) (5.239,5.024) (5.297,4.974) (5.355,4.924) (5.411,4.872) (5.467,4.820) (5.522,4.767) (5.576,4.712) (5.629,4.657) (5.680,4.600) (5.730,4.541) (5.777,4.481) (5.823,4.420) (5.866,4.357) (5.907,4.292) (5.946,4.226) (5.982,4.158) (6.015,4.088) (6.045,4.017) (6.073,3.944) (6.098,3.870) (6.119,3.794) (6.140,3.719) (6.165,3.646) (6.193,3.576) (6.226,3.508) (6.262,3.442) (6.302,3.379) (6.346,3.318) (6.394,3.258) (6.446,3.200) (6.502,3.143) (6.561,3.087) (6.623,3.032) (6.685,2.976) (6.745,2.919) (6.803,2.860) (6.859,2.799) (6.914,2.737) (6.967,2.674) (7.019,2.610) (7.069,2.545) (7.119,2.479) (7.167,2.412) (7.213,2.344) (7.259,2.275) (7.303,2.205) (7.345,2.134) (7.386,2.062) (7.426,1.990) (7.464,1.916) (7.500,1.842) (7.535,1.767) (7.568,1.692) (7.598,1.615) (7.627,1.538) (7.654,1.460) (7.678,1.382) (7.701,1.302) (7.721,1.223) (7.739,1.143) (7.756,1.062) (7.770,0.982) (7.783,0.901) (7.794,0.819) (7.804,0.738) (7.812,0.656) (7.819,0.574) (7.825,0.492) (7.830,0.410) (7.834,0.328) (7.837,0.246) (7.839,0.164) (7.842,0.082) (7.843,0.000)};}
    \capa{12}{\node[etiqueta,text=tceeazul,anchor=south west] at (6.1,4.35) {GFPU};}
  \end{scope}
  % E_1 y E_3 en la FPP, con su FPU; E es E_2
  \capa{9}{\node[punto,label={[etiqueta]above left:$E_1$}] at (3,5.421047) {};}
  \capa{10}{\node[etiqueta,anchor=base west,inner xsep=0pt] at (Elab.base east) {${}=E_2$};}
  \capa{11}{\node[punto,label={[etiqueta]above right:$E_3$}] at (6.3,2.615339) {};}
\end{tikzpicture}
\end{document}
